% F8: Research plan. 1.5 to 2 pages. Winter (W4): the basis of your methods section.
% Checklist: testable questions or hypotheses; a row in the evaluation table
% for every main claim; at least one strong, cited baseline; at least 3 risks
% with mitigations; dated milestones; code and data plan with repository
% link; ethics and data-protection note; AI critique recorded.
% Record the AI critique ("What would make this evaluation unconvincing to a
% reviewer?") and your changes in appendices/assignment_records.tex.

\section{Research Plan}

\subsection{Research Questions and Approach}
\textit{[F8: research questions or hypotheses, and your approach.]}

\subsection{Evaluation}
\tagpdfsetup{table/header-rows={1}}
\begin{longtable}{@{}p{0.22\linewidth}p{0.2\linewidth}p{0.14\linewidth}p{0.18\linewidth}p{0.18\linewidth}@{}}
  \toprule
  Claim & Experiment & Metric & Baseline & Dataset or benchmark \\
  \midrule
  \endhead
  \textit{[claim]} & \textit{[experiment]} & \textit{[metric]} & \textit{[baseline, cited]} & \textit{[dataset]} \\
   & &&& \\
  \bottomrule
\end{longtable}

\subsection{Risks}
\begin{itemize}
  \item \textit{[Risk]}: \textit{[mitigation]}
\end{itemize}

\subsection{Milestones}
\textit{[F8: dated milestones up to your RPE or thesis defense.]}

\subsection{Code and Data}
\textit{[F8: repository link and structure, environment, data storage, anything confidential.]}
\tagpdfsetup{table/header-rows={1}}
\begin{longtable}{@{}p{0.22\linewidth}p{0.2\linewidth}p{0.14\linewidth}p{0.18\linewidth}p{0.18\linewidth}@{}}
  \toprule
  Code Storage & Repo link\\
  \midrule
  \endhead
    Particle Filter and UKF experiments & \href{https://github.com/camstillo/ParticleFilterAttitudeEstimator.git}{Github Link} \\
    Power scheduler simulation & \href{https://github.com/camstillo/PowerSchedulerAlgorithms.git}{Github Link} \\
    Pointing Optimization Code & \href{https://github.com/camstillo/PointingOptimizationCubesat.git}{Github Link} \\
  \bottomrule
\end{longtable}

Currently, these don't have a directory structure, as the code hasn't necessarily been started on. These will eventually have a directory structure located in the README, but I usually tend to have a single folder with all supporting code (includes, etc.), a build folder for compilation, and a \texttt{main.c} file at the top level for compilation. Any dependencies will be consistently managed generally with a single version.
 
For dependencies, I am planning on using IPOPT with MUMPS and BLAS. Versions of these will likely not change at all from the version that I'm currently using. I plan on using GNU Scientific Library (GSL) to handle some of the linear algebra stuff that I'll be doing with the code. Finally, \href{https://www.space-track.org/}{SpaceTrack.org} gives out an SGP4 orbit propagator for use with a C interface, which is what I plan on using for the orbital data simulator (version shouldn't change).

\subsection{Ethics and Data Protection}
\textit{[F8: 2 to 3 sentences: IRB status if you work with human subjects; privacy and export-control constraints; who may see your data.]}

I am not going to be working with human subjects or anything involving privacy. None of the work I'm doing is export controllable, although I need to ensure to not include the SGP4 orbital propagator in any git repositories as that is not available 
